% !TeX root = ../main.tex
\chapter{Related Work}
\label{ch:relatedwork}
Die Automatisierung komplexer Software-Migrationen erfordert ein tiefgreifendes Verständnis sowohl der proprietären Zielsysteme (hier SAP) als auch aktueller KI-Technologien. Da die direkte Verbindung von SAP-Migrationen und generativer KI bislang wenig erforscht ist, ordnet dieses Kapitel die vorliegende Arbeit anhand aktueller Publikationen aus verwandten Disziplinen in den Stand der Technik ein. Zunächst analysiert das Kapitel die konkreten Hürden bei der Transformation von SAP-Konfigurationsmodellen und beleuchtet die Limitierungen aktueller Standardwerkzeuge. Darauf aufbauend wird der Einsatz von Large Language Models (LLMs) für die Code-Migration betrachtet. Dabei wird aufgezeigt, warum monolithische Ansätze bei strengen Architekturvorgaben in der Praxis häufig scheitern. Abschließend stellt das Kapitel den aktuellen Forschungsstand zu Multi-Agenten-Systemen (MAS) vor und leitet daraus die zentrale Forschungslücke dieser Arbeit ab: die semantische Migration proprietärer SAP-Constraint-Logik.

Bei der Migration auf SAP S/4HANA stehen Unternehmen vor der Aufgabe, ihre historisch gewachsenen Produktmodelle von der klassischen Variantenkonfiguration (LO-VC) in die Advanced Variant Configuration (AVC) überführen. \textcite[S. 31]{matilainenChangeRequirementsProduct2024} zeigt in seiner Studie die Komplexität dieses Wechsels auf: Da die Zahl der möglichen Konfigurationen exponentiell mit den Merkmalen und Optionen wächst, ist eine präzise Restrukturierung der Datenmodelle erforderlich, um die Variantenvielfalt im neuen System technisch beherrschen zu können.

Ein wesentlicher Schritt vor der eigentlichen Migration ist die Datenbereinigung. Laut \textcite[S. 52, S. 84]{matilainenChangeRequirementsProduct2024} müssen ungenutzte Altdaten (Legacy Data) zwingend bereinigt werden, um Strukturen zu vereinfachen und das Zielsystem zu entlasten. Obgleich SAP für diesen Umstieg Standardwerkzeuge wie Transition Workspaces oder das Migration Cockpit bereitstellt, liegt deren Fokus primär auf der technischen Übernahme von Massendaten. Das in Kapitel 3 geführte Experteninterview identifiziert hier eine signifikante Lücke: Die SAP-Werkzeuge arbeiten weitgehend regelbasiert und deterministisch. Sie sind nicht in der Lage, das historische, prozedurale Beziehungswissen autonom zu interpretieren und in die neue, deklarative AVC-Logik zu übersetzen. Diese inhaltliche Restrukturierung sowie die Bereinigung von Legacy-Code erfordern derzeit manuelle, ressourcenintensive Eingriffe, für die in der industriellen Praxis bislang keine automatisierten Lösungen existieren.

Die Entwicklungen im Bereich der Large Language Models (LLMs) haben die Möglichkeiten zur automatisierten Code-Generierung und -Transformation grundlegend erweitert \parencite{karabiyikRefactorGPTChatGPTbasedMultiagent2025}. Die aktuelle Forschung belegt jedoch, dass die Migration komplexer Altsysteme nicht durch eine einzige LLM-Abfrage (Single-Prompt oder Monolithic Prompting) bewältigt werden kann. Derartige Ansätze scheitern in der Praxis häufig an der unzureichenden Berücksichtigung spezifischer Architekturvorgaben und systeminterner Abhängigkeiten. \textcite{motiLegacyTranslateLLMbasedMultiAgent2026} demonstrieren am Beispiel des Frameworks LegacyTranslate, dass naive LLM-Übersetzungen bei industriellen Codebasen Kompilierraten von 0 \% aufweisen können. Obwohl der generierte Code syntaktisch plausibel erscheint, ist er aufgrund fehlender interner Schnittstellen (APIs) und domänenspezifischer Logik nicht funktionsfähig \parencite{motiLegacyTranslateLLMbasedMultiAgent2026}. Ähnliche Defizite dokumentieren \textcite{xuMANTRAEnhancingAutomated2025}: Während spezialisierte Systeme bei strukturellen Code-Änderungen hohe Erfolgsquoten erzielen, erreichen monolithische Basismodelle lediglich eine Erfolgsquote von 8,7 \%.

Ein weiteres konzeptionelles Problem dieser One-Shot-Ansätze liegt in ihrer mangelnden Kontrollierbarkeit \parencite{karabiyikRefactorGPTChatGPTbasedMultiagent2025, oueslatiRefAgentMultiagentLLMbased2026}. Ohne eine Modularisierung in kleinere Arbeitsschritte bleibt der Transformationsprozess eine intransparente Blackbox. Es können keine funktionalen Garantien für den Zielcode abgegeben werden, und die Modelle neigen bei komplexeren Interaktionen zu Halluzinationen \parencite{oueslatiRefAgentMultiagentLLMbased2026,karabiyikRefactorGPTChatGPTbasedMultiagent2025}. Um diese Unzulänglichkeiten zu überwinden, fokussiert sich die aktuelle Forschung zunehmend auf Multi-Agenten-Systeme (MAS). Deren Grundprinzip besteht darin, einen komplexen Prozess in logische Teilaufgaben zu zerlegen, die anschließend von kooperierenden, spezialisierten KI-Agenten autonom bearbeitet werden.

Erste Studien im Bereich des automatisierten Software-Refactorings – der strukturellen Optimierung innerhalb derselben Programmiersprache – belegen die prozessuale Überlegenheit dieses modularen Ansatzes. Das Framework RefactorGPT überführt den Refactoring-Prozess durch eine sequentielle Orchestrierung in einen kontrollierbaren Workflow \parencite{karabiyikRefactorGPTChatGPTbasedMultiagent2025}. Auch das Framework RefAgent validiert diesen Vorteil: Die Aufteilung in spezialisierte Agenten für Planung, Ausführung und Fehlerbehebung steigert die Kompiliererfolgsrate signifikant und reduziert fehlerhafte Generierungen \parencite{oueslatiRefAgentMultiagentLLMbased2026}. Ein zentraler Erfolgsfaktor sind dabei autonome Feedback-Schleifen. So integriert das System MANTRA einen „Repair Agent“, der Compiler-Fehler iterativ korrigiert, wodurch die Erfolgsquote bei komplexen Transformationen auf 82,8 \% gesteigert werden konnte \parencite{xuMANTRAEnhancingAutomated2025}. Obwohl ein solches direktes Compiler-Feedback in der vorliegenden Arbeit aufgrund restriktiver, proprietärer SAP-Systemarchitekturen noch der Future Work vorbehalten bleibt, unterstreichen diese Ergebnisse das technologische Potenzial von MAS.

Während die Refactoring-Ansätze die Effizienz der Agenten-Kollaboration validieren, demonstriert das Projekt LegacyTranslate \parencite{motiLegacyTranslateLLMbasedMultiAgent2026}, dass sich diese Architektur auch auf deutlich komplexere Cross-Language-Migrationen übertragen lässt. Die Übersetzung von PL/SQL zu Java erfordert hier einen Paradigmenwechsel. LegacyTranslate implementiert hierfür ein „API Grounding“, durch welches das System aktiv an die Architekturvorgaben des Zielsystems gebunden wird. 

Trotz der empirisch belegten Leistungsfähigkeit von MAS bei der Code-Migration offenbart die aktuelle Literatur eine signifikante Lücke hinsichtlich proprietärer ERP-Ökosysteme. Etablierte Frameworks wie MANTRA oder LegacyTranslate fokussieren sich primär auf imperative oder objektorientierte Standardsprachen (wie Java oder PL/SQL). Bislang wurde nicht wissenschaftlich untersucht, inwiefern kollaborative KI-Architekturen auf die domänenspezifische Constraint-Logik der SAP angewendet werden können. Es bleibt eine offene Forschungsfrage, ob Multi-Agenten-Systeme fähig sind, historisch gewachsene Regelwerke semantisch in ein deklaratives Format – ein Constraint Satisfaction Problem (CSP) – zu überführen. Für diesen notwendigen Paradigmenwechsel, der LO-VC-Vorbedingungen in ein mathematisches Lösungsnetzwerk transformiert, bieten bestehende KI-Ansätze noch keine architektonische Lösung.

An dieser methodischen Schnittstelle setzt die vorliegende Arbeit an. Sie adressiert die identifizierte Forschungslücke, indem sie das Konzept verteilter Agenten-Architekturen auf eine spezifische Herausforderung in der Migration der SAP-Variantenkonfiguration überträgt. 